\section*{Chapitre \ref{ch:ml:alternative-data-pipelines} --- Chaînes d'ingestion des données alternatives}

\begin{solution}{exo:ml:alternative-data-pipelines:1}
Le ratio vaut environ 100, bien en dehors de $[0.5, 2]$~: la livraison est retenue pour examen. Un facteur de 100 indique un changement
d'unités (centimes pour des dollars)~; les autres causes d'un ratio élevé sont un fichier livré deux fois, des données d'un nouveau partenaire ajoutées, ou
des totaux cumulés envoyés au lieu des accroissements.
\end{solution}

\begin{solution}{exo:ml:alternative-data-pipelines:2}
Pondérations $\begin{pmatrix}1.18 & 0.72\\ 1.45 & 0.89\end{pmatrix}$ (à deux décimales près), effectifs pondérés
$\begin{pmatrix}35.5 & 14.5\\ 14.5 & 35.5\end{pmatrix}$~: chaque ligne et chaque colonne totalisent 50. Le raking ajuste les marges, et non
les cellules~; beaucoup de pondérations y parviennent, et le raking retient celle qui est la plus proche des pondérations de départ au sens de l'entropie relative.
\end{solution}

\begin{solution}{exo:ml:alternative-data-pipelines:3}
«~BOREALIS 5120~» se normalise en «~borealis~», qui correspond de la même façon au premier jeton du nom de trois distributeurs (Borealis
Retail, Stores and Outfitters) et aucun de leurs seconds jetons~: score de 0.79 pour chacun, une égalité dans la zone de revue. «~BOREALI RETA
STR 313~» se normalise en «~boreali reta~», dont les jetons sont des préfixes de «~borealis~» et de «~retail~»~: score de 1 pour Borealis
Retail seul.
\end{solution}

\begin{solution}{exo:ml:alternative-data-pipelines:4}
La croissance d'une année sur l'autre compare un trimestre au même trimestre un an plus tôt. Pour les trimestres 16 à 19, la valeur courante provient
du nouveau panel et la valeur antérieure de l'ancien, de sorte que le changement de composition du panel est lu comme une variation des ventes, dans
une mesure qui dépend de la clientèle de chaque distributeur. À partir du trimestre 20, les deux termes proviennent du nouveau panel, et l'erreur
revient au niveau d'un panel stable (2.7~points).
\end{solution}

\begin{solution}{exo:ml:alternative-data-pipelines:5}
Il a été livré au lancement, un trimestre après sa fin, exactement comme un trimestre en temps réel, si bien que ses horodatages sont indiscernables
de ceux des données en temps réel. Cela n'a pas eu d'importance, car le fournisseur ne pouvait pas l'ajuster~: les distributeurs n'avaient pas encore déclaré leurs ventes. Cela aurait
de l'importance si le fournisseur l'avait construit avec une information postérieure (une correspondance des marchands ajustée après coup, un panel nettoyé des
\omterm{def:ml:neural-networks-for-noisy-tabular-data:dropout}{dropouts} postérieurs), ce que seule la documentation du fournisseur ou une comparaison des millésimes peut révéler.
\end{solution}

\begin{solution}{exo:ml:alternative-data-pipelines:6}
Réparer masque le problème du fournisseur et devine la correction~: une dépense négative peut être un remboursement, une erreur de signe ou une
ligne corrompue, et sa valeur absolue n'en est qu'une parmi d'autres. Rejeter, compter par motif, signaler au fournisseur, et ne réparer qu'avec une règle
confirmée par le fournisseur, consignée avec les données.
\end{solution}

\begin{solution}{exo:ml:alternative-data-pipelines:7}
L'erreur après raking passe de 2.1 à 19.7~points en moyenne. Les chaînes écartées portaient 18.4~\% des dépenses livrées
par les distributeurs, mais dans une part qui change d'un trimestre à l'autre (la dépense de chaque cellule est répartie au hasard entre deux ou trois
chaînes d'un même distributeur), si bien que les écarter ajoute du bruit à chaque taux de croissance plutôt qu'un biais constant. Les distributeurs et les trimestres
pour lesquels la chaîne réduite au seul radical portait par hasard davantage de dépenses sont les plus touchés.
\end{solution}

\begin{solution}{exo:ml:alternative-data-pipelines:8}
Soient $n_{ar} = N u_a v_r$ les effectifs des cellules du panel et $p_{ar} = \alpha_a\beta_r$ les parts de la population. Les pondérations
$w_{ar} = (\alpha_a/u_a)(\beta_r/v_r)$ donnent des effectifs pondérés $N\alpha_a\beta_r$, qui respectent les deux marges, et les deux premières étapes du raking y parviennent~:
l'étape sur les lignes multiplie par $\alpha_a/u_a$, après quoi les sommes par colonne valent $N\alpha_\cdot v_r$ et la
deuxième étape, sur les colonnes, multiplie par $\beta_r/v_r$. Ce sont les pondérations de post-stratification $p_{ar}/(n_{ar}/N)$. Lorsque la population
présente une interaction ($p_{ar}\ne\alpha_a\beta_r$), le raking respecte toujours les marges mais pas les cellules, et une statistique qui
dépend de l'interaction reste biaisée.
\end{solution}

\begin{solution}{pb:ml:alternative-data-pipelines:1}
\textbf{Partie I.}
\begin{enumerate}
\item Champs manquants, dépenses négatives, effectifs du panel saisis en texte, lignes dupliquées et une livraison en centimes. La validation rejette
121, 69, 64 et 111 lignes, soit chacun des défauts injectés.
\item La livraison en centimes~: toutes les lignes sont valides. Le contrôle de livraison, comparé au total précédent, la retient.
\item Une réparation est une supposition sur l'erreur du fournisseur~; le rejet préserve la fiabilité des données et oblige le fournisseur à la corriger.
\item La valeur, avec son temps de validité, et l'heure de livraison comme temps de connaissance, de sorte que chaque requête puisse être posée à une date donnée
et que le rétropeuplement soit visible.
\end{enumerate}
\textbf{Partie II.}
\begin{enumerate}[resume]
\item Appariement exact sur le nom normalisé~: 300 sur 600. Le résolveur~: 500 appariés automatiquement, tous correctement, 100 à
examiner, aucun sans correspondance.
\item Les chaînes qui n'ont en commun avec jusqu'à trois distributeurs qu'un premier mot~: une machine ne peut pas les distinguer, mais une personne
munie des listes de magasins le peut.
\item Aucune n'est appariée~; six sont renvoyées en revue, quatre sont rejetées.
\item Un mauvais appariement déplace silencieusement des dépenses d'une entreprise à l'autre et biaise les deux nowcasts~; une chaîne en revue retarde les données et
coûte du temps à une personne.
\end{enumerate}
\textbf{Partie III.}
\begin{enumerate}[resume]
\item Brut~: 3.3~points (trimestres 4--15), 22.5 (16--19), 2.7 (20--23). Après raking~: 1.7, 2.7, 2.6.
\item Que les marges de la population sont connues, que le panel couvre chaque cellule, et qu'à l'intérieur d'une cellule les panélistes dépensent comme
la population. Le résidu tient au bruit d'échantillonnage dans les petites cellules et à toute différence au sein d'une cellule.
\item Trimestres 0 à 10~: leur \omterm{def:ml:alternative-data-pipelines:firstseen}{horodatage de première observation} tombe au trimestre 12, soit plus d'un trimestre après leur fin.
\item Après raking~: 0.480 rétropeuplé, 0.419 en temps réel avant le changement, 0.412 après. Brut~: 0.366, 0.391, 0.264.
\end{enumerate}
\textbf{Partie IV.}
\begin{enumerate}[resume]
\item \emph{Le panel qui a grossi.} Dans l'année qui suit le changement de partenaire, le nowcast brut manque la croissance de 22.5~points en
moyenne, celui après raking de 2.7~points~; l'IC de l'historique rétropeuplé vaut 0.480, contre 0.419 en temps réel.
\item La composition du panel par période et les marges auxquelles il est pondéré~; les changements de partenaire et de méthode, avec leurs dates~;
les horodatages de livraison de tout l'historique~; la carte des marchands et ses évolutions.
\item Une rupture de composition (effectifs du panel par cellule), un écart dans le lien entre les totaux du panel et les ventes déclarées (CUSUM du livre~7 sur
les résidus), et un saut dans la part des nouvelles chaînes de marchands.
\item Les contrôles de livraison, les taux de rejet, la file de revue, la composition du panel, les résidus du nowcast par rapport aux ventes déclarées, et
l'IC en temps réel.
\item L'écart avec le temps réel n'est en rien négligeable~; il faut traiter l'historique rétropeuplé comme une borne supérieure et tabler sur le chiffre en temps réel.
\item Sur la file de revue et les livraisons retenues~: là où l'incertitude de la machine est concentrée et où le savoir d'une personne
est le moins coûteux.
\item Les mêmes étapes avec d'autres entités (sites, navires, applications) et d'autres biais (couverture nuageuse, répartition des appareils, ouvertures de
magasins)~; la question de la composition et celle de la première observation ne disparaissent pas.
\item Que le panel représente encore la population qu'il représentait lorsque le modèle a été ajusté.
\end{enumerate}
\end{solution}

\begin{solution}{iq:ml:alternative-data-pipelines:1}
Historique livré après coup, souvent construit avec une information indisponible à l'\omterm{def:ml:neural-networks-for-noisy-tabular-data:backprop}{époque} (une carte des marchands, calibration sur
les chiffres déclarés). On le détecte en enregistrant les \omterm{def:ml:alternative-data-pipelines:firstseen}{horodatages de première observation} à l'ingestion, en demandant au fournisseur les dates de livraison et
les millésimes, et en comparant les statistiques de l'historique à celles de la période en temps réel.
\iqlookfor{la définition, les \omterm{def:ml:alternative-data-pipelines:firstseen}{horodatages de première observation} et la comparaison entre temps réel et rétropeuplement.}
\end{solution}

\begin{solution}{iq:ml:alternative-data-pipelines:2}
\omterm{def:ml:alternative-data-pipelines:ingest}{Validation de schéma} (champs, types, plages de valeurs, unicité), contrôles au niveau de la livraison (nombre de lignes, totaux comparés à la dernière livraison, couverture, heure d'arrivée),
et contrôles de référence (entités inconnues). En cas d'échec~: rejeter les lignes en indiquant les motifs, suspendre la livraison, alerter,
ne jamais réparer en silence~; stocker avec les temps de connaissance.
\iqlookfor{contrôles sur les lignes et sur les livraisons, et politique d'échec sans réparation silencieuse.}
\end{solution}

\begin{solution}{iq:ml:alternative-data-pipelines:3}
Normaliser les descripteurs, bloquer les candidats par jetons, noter par similarité de jetons et de préfixes (ou par un appariement entraîné), accepter
au-dessus d'un seuil, soumettre à révision une zone intermédiaire, rejeter en dessous~; conserver la correspondance à chaque instant.
Mesurer la précision et le rappel sur un échantillon étiqueté, stratifié par motif de descripteur, et suivre la file de révision et les taux de nouvelles chaînes.
\iqlookfor{seuils avec révision, correspondance à chaque instant, précision mesurée.}
\end{solution}

\begin{solution}{iq:ml:alternative-data-pipelines:4}
Repondérer le panel sur les marges de population (calage sur marges ou post-stratification) par âge et selon d'autres dimensions connues. La méthode échoue
lorsque les cellules sont vides ou très petites, lorsque les marges sont inconnues ou périmées, lorsque les panélistes diffèrent de la population au sein des cellules, et
lorsque les changements de composition rendent caduques les comparaisons d'une année sur l'autre.
\iqlookfor{repondération et ses hypothèses, y compris les écarts au sein des cellules.}
\end{solution}

\begin{solution}{iq:ml:alternative-data-pipelines:5}
Les données avant le modèle~: un changement de partenaire ou de méthodologie, la composition du panel, la carte des marchands, le calendrier des livraisons~; puis le distributeur
(un nouveau canal que le panel ne voit pas, comme les ventes en ligne passant par un autre processeur)~; seulement ensuite le modèle.
\iqlookfor{changements de panel et de correspondance avant les changements de modèle.}
\end{solution}

\begin{solution}{iq:ml:alternative-data-pipelines:6}
Scinder l'essai au lancement~: les deux années en temps réel constituent la preuve, avec leur erreur type~; les six années rétropeuplées forment une borne supérieure.
Vérifier la méthode de rétropeuplement du fournisseur, comparer les millésimes, et ramener l'estimation vers l'IC en temps réel~; s'attendre
à une nouvelle érosion à mesure que d'autres acquièrent les données (livre~8, chapitre~16).
\iqlookfor{estimation sur le seul temps réel, scepticisme envers le rétropeuplement, et dégradation attendue.}
\end{solution}
